\documentclass[10pt,a4paper]{amsart}
\usepackage[margin=19mm]{geometry}
\usepackage{amsmath,amssymb,mathtools}
\usepackage[scr=rsfs,cal=euler]{mathalfa}
\usepackage{xfrac}
\usepackage[unicode,hidelinks,bookmarks=false]{hyperref}
\newcommand{\ie}{i.e.\ }
\newcommand{\cf}{cf.\ }
\newcommand{\resp}{respectively\ }
\newcommand{\Subsection}{\subsection}
\providecommand{\nobreakdash}{\nobreak\hbox{-}}
\setlength{\emergencystretch}{2em}
\multlinegap0pt
\usepackage{bm}
\usepackage{bbm}
\usepackage{dsfont}
\usepackage{xspace}
\usepackage{cancel}
\usepackage{mathtools}
\usepackage{amsthm}
\makeatletter
\@ifundefined{theorem}{\newtheorem{theorem}{Theorem}}{}
\@ifundefined{lemma}{\newtheorem{lemma}[theorem]{Lemma}}{}
\makeatother
\makeatletter
\def\E{{\mathbb E}}
\def\R{{\mathbb R}}
\@ifundefined{assumption}{\newenvironment{assumption}[1]
 {%
  \renewcommand{\assumptionnumber}{Assumption #1}%
  \begin{assumption*}%
  \protected@edef\@currentlabel{\textbf{\textup {#1}}}%
 }}{}
\makeatother
\title[Sharp propagation of chaos for mean field Langevin dynamics,]{Sharp propagation of chaos for mean field Langevin dynamics, control, and games}
\author{Manuel Arnese, Daniel Lacker}
\date{Source: arXiv:2603.10988v1 (CC BY 4.0)}
\begin{document}
\maketitle

\noindent\emph{Source and licence.} Sharp propagation of chaos for mean field Langevin dynamics, control, and games. Version arXiv:2603.10988v1 (CC BY 4.0), distributed under the Creative Commons Attribution 4.0 International licence, as recorded in the arXiv licence metadata for this exact version and shown on the abstract page https://arxiv.org/abs/2603.10988v1. Full rights notice: licenses/arxiv-2603.10988-CC-BY-4.0.txt.\par\smallskip

\noindent\textit{Editorial scope.} Counted unit: the Lemma whose source label is \texttt{lm: linear differential inequality} in the article (source0.tex lines 655--666, source label \texttt{lm: linear differential inequality}), delivered together with the article's own complete original proof of it (source0.tex lines 667--726, one \texttt{proof} environment of 60 lines). No mathematics is written, repaired or re-derived: statement and proof are the article's own text line for line. Required context retained uncounted: eq: linear differential inequality (lines 649--652). Every definition, display and label of those lines is retained unchanged. Omitted from this package and not redistributed: 1--648, 653--654, 727--1856. Nothing inside a retained excerpt is omitted or altered: every retained body is delivered exactly as the article's lines are written, so no omission receipt of any kind accompanies this package. Citations: the retained excerpts carry no citation command at all, so no bibliography entry of the article is transcribed and no reference is redistributed. Reference adaptation: none was needed and none was made; every reference inside the delivered unit resolves inside it, so no unresolved reference or citation is produced. Printed slips kept literal: no printed word, symbol, hypothesis or number of the counted statement or of its proof was altered. Layout and class: the article's own class is \texttt{amsart} with packages including amsmath, amsthm, graphicx, amsfonts, amssymb, hyperref, bm, color, geometry, enumerate, mathrsfs, natbib, mathtools; the wrapper is the lane's pinned \texttt{amsart} editorial wrapper at 10pt on A4, which restates the theorem-like environments the retained text needs and adds the article's own macro definitions that the retained text needs (\texttt{\textbackslash E}, \texttt{\textbackslash R}), with extra wrapper packages (bm, bbm, dsfont, xspace, cancel, mathtools). The delivered numbering is the unit's own. Assets: no retained span contains a figure environment, a graphics-inclusion command, a PDF-page inclusion, a table or a TikZ picture; no image file is copied into this package, the package declares no asset dependency and no image file is redistributed with it. Licence: version arXiv:2603.10988v1 of this article is distributed under the Creative Commons Attribution 4.0 International licence, as recorded in the arXiv licence metadata for this exact version and shown on the abstract page https://arxiv.org/abs/2603.10988v1. A full rights notice is delivered at licenses/arxiv-2603.10988-CC-BY-4.0.txt. Characters: every retained excerpt is pure ASCII, so the delivered engine, XeLaTeX with its default Latin Modern text font, reports no missing character.
\begin{equation}
\label{eq: linear differential inequality}
    \frac{d}{dt}f_t(k)\leq ak(f_t(k+1)-f_t(k))1_{k < n}  + \frac{b}{n^2}k^p + kR_t - cf_t(k).
\end{equation}

\begin{lemma}
  \label{lm: linear differential inequality}
  Assume that $f:[0,T]\times [n]\to \R_+$ is a non-decreasing function in $k$ that satisfies \eqref{eq: linear differential inequality} for $p\in\{2,3\}$. Assume that there exists a constant $C_0\geq 0$ such that 
  \begin{align}
  \label{eq: yule process initial condition}
      f_s(k)\leq C_0\frac{k^2}{n^2}, \quad 1 \le k \le n.
  \end{align}
  Then, for $0 \le s \le T$ and $1 \le k \le n$
    \begin{equation*}
        f_T(k)\leq 2C_0\frac{k^2}{n^2} e^{(2a-c)(T-s)}+\frac{8bk^p}{n^2}\int_s^T e^{(ap-c)(t-s)} \,dt + k\int_s^T e^{(a-c)(t-s)} R_t \,dt.
    \end{equation*}
\end{lemma}

\begin{proof}
We will use a semigroup interpretation of \eqref{eq: linear differential inequality} to organize the analysis.
Consider the $n \times n$ matrix $\mathcal G$ which acts on a function $\varphi : [n] \to \R$ by
\[\mathcal G \varphi(k)=ak(\varphi(k+1)-\varphi(k))\mathbf{1}_{k<n}.\]
Each row of this matrix sums to zero (i.e., $\mathcal G$ maps constant functions to zero), and every off-diagonal entry is nonnegative. Hence, $\mathcal G$ is the infinitesimal generator of a continuous-time Markov process on state space $[n]$; in fact, this process goes by the name of the \emph{Yule process} stopped at level $n$, and we denote it by $(\mathcal Y_t)_{t\geq 0}$. The matrix exponential is the semigroup of this process and admits the representation $e^{t\mathcal G}\varphi(k)=\E[\varphi(\mathcal Y_t)\,|\,\mathcal Y_0=k]$. We will not use this stochastic interpretation except to deduce a monotonicity property, namely that $e^{t\mathcal G}\varphi \ge e^{t\mathcal G}\psi$ pointwise whenever $\varphi \ge \psi$ pointwise.
With this notation we can rewrite \eqref{eq: linear differential inequality} as a coordinatewise inequality between vectors,
\begin{equation*}
    \frac{d}{dt}f_t \leq (\mathcal G - cI)f_t + \frac{b}{n^2}m_p + R_tm_1,
\end{equation*}
where $I$ denotes the $n \times n$ identity matrix, and $m_q(k) := k^q$ for $q > 0$.
Using the aforementioned monotonicity, we deduce
\begin{align*}
    \frac{d}{dt}\Big[e^{-ct} e^{t\mathcal G}f_{T-t} \Big] &= e^{-ct} e^{t\mathcal G}\Big[(\mathcal G - cI) f_{T-t}-\frac{d}{ds}f_{s}\big\vert_{s=T-t}\Big] \\
    &\ge - e^{-ct}e^{t\mathcal G}\big(\frac{b}{n^2}m_p + R_tm_1\big).
\end{align*}
Integrate to get the coordinatewise inequality
\begin{align}
\label{eq: semigroup after gronwall}
    f_{T} &\le e^{-c(T-s)}e^{(T-s)\mathcal G} f_s + \int_s^Te^{-c(t-s)}e^{(t-s)\mathcal G}\Big(\frac{b}{n^2} m_p + R_tm_1\Big)\,dt.
\end{align}
Now we estimate $e^{t\mathcal G}m_q$ for $q=1,2,3$. This could be done by using the known fact that the Yule process follows a negative binomial distribution, but we instead present a short self-contained argument. 

\begin{itemize}
\item $q=1$: Use $\mathcal G m_1(k) = ak(m_1(k+1)-m_1(k))=ak$ to get
\begin{align*}
\frac{d}{dt}e^{t\mathcal G}m_1=e^{t\mathcal G}\mathcal Gm_1 = ae^{t\mathcal G}m_1,
\end{align*}
and thus $e^{t\mathcal G}m_1=e^{at}m_1$.
\item $q=2$: We have $\mathcal G m_2(k) = ak\big((k+1)^2-k^2\big) = 2a m_2(k) + am_1(k)$, and thus
\begin{align*}
\frac{d}{dt}e^{t\mathcal G}m_2 &= 2ae^{t\mathcal G}m_2 + a e^{t\mathcal G}m_1.
\end{align*}
Multiply by $e^{-2at}$ and integrate to get
\begin{align*}
e^{t\mathcal G}m_2 &= e^{ 2at}m_2 + ae^{2at}\int_0^t e^{-2as} e^{s\mathcal G}m_1 \,ds.
\end{align*}
Using $e^{t\mathcal G}m_1=e^{at}m_1$, we get the coordinatewise inequality
\[
e^{t\mathcal G}m_2 = e^{ 2at}m_2 + e^{2at}(1-e^{-at}) m_1 \le 2 e^{2at}m_2.
\]
\item $q=3$: We have $\mathcal G m_3(k) = ak\big((k+1)^3-k^3\big) = 3a m_3(k) + 3am_2(k) + am_1(k)$, and thus
\begin{align*}
\frac{d}{dt}e^{t\mathcal G}m_3 &= 3ae^{t\mathcal G}m_3 + 3a e^{t\mathcal G}m_2 + a e^{t\mathcal G}m_1.
\end{align*}
Multiply by $e^{-3at}$ and integrate to get
\begin{align*}
e^{t\mathcal G}m_3 &= e^{ 3at}m_3 + ae^{3at}\int_0^t e^{-3as} \big(3  e^{s\mathcal G}m_2 +   e^{s\mathcal G}m_1\big) \,ds.
\end{align*}
Use $e^{s\mathcal G}m_1\le e^{2as} m_3$ and $e^{s\mathcal G}m_2 \le 2 e^{2as}m_3$ to get
\begin{align*}
e^{t\mathcal G}m_3 &\le \bigg[e^{ 3at} + 7ae^{3at}\int_0^t e^{-as} \,ds\bigg] m_3 \le 8 e^{3at} m_3.
\end{align*}
\end{itemize}
Using the $q=2$ case and the time-$s$ assumption, we have $e^{(T-s)\mathcal G} f_s(k) \le 2C_0 e^{2a(T-s)}k^2/n^2$.
We can summarize all cases $q=1,2,3$ as $e^{t\mathcal G}m_q \le 8e^{qat}m_q$.
Combining these bounds with \eqref{eq: semigroup after gronwall} yields
\begin{align*}
f_{T}(k) \le 2C_0 e^{(2a-c)(T-s)} \frac{k^2}{n^2}  + \int_s^Te^{-c(t-s)}\Big(\frac{8bk^p}{n^2}e^{pa(t-s)}   + R_tke^{a(t-s)}\Big)\,dt.
\end{align*}
\end{proof}

\end{document}
